\chapter{Resumé}
\label{ch:abstract_da}

Systemer med interagerende komponenter modelleres ofte som netværk, både inden for naturvidenskabelige felter såsom biologi og fysik og inden for samfundsvidenskab. Transportprocesser i disse netværk er af særlig interesse og undersøges ofte ved hjælp af flow-optimering. Denne afhandling fokuserer på flow i netværk bestående af kemiske reaktioner - såkaldte \emph{reaktionsnetværk} - hvor de enkelte molekylers struktur er kendt. Molekylernes struktur har betydning for, hvilke reaktionsveje der kan identificeres i netværket. Reaktionsnetværk kan naturligt repræsenteres som hypergrafer, hvilket giver et matematisk grundlag, som understøtter brugen af computere til at finde reaktionsveje.

Reaktionsnetværks størrelse, struktur og kombinatoriske kompleksitet gør dem vanskelige at analysere manuelt. Et centralt problem er at identificere mulige reaktionsveje fra en given mængde af udgangsmolekyler til et bestemt målmolekyle, under overholdelse af de underliggende fysiske og kemiske love og begrænsninger. At finde reaktionsveje, der optimerer produktionen af et målmolekyle ud fra enten termodynamiske hensyn eller reaktionssandsynligheder, er et grundlæggende problem i mange sammenhænge, herunder i kemiske reaktorsystemer. Denne afhandling giver tre metodiske bidrag. For det første udvides eksisterende teknikker for søgning efter reaktionsveje til at inddrage termodynamisk information. For det andet integreres kinetisk information i en rangordningen af de fundne reaktionsveje. For det tredje introduceres en stokastisk metode til udforskning af implicit definerede reaktionsnetværk.

De foreslåede metoder kombinerer veletablerede værktøjer fra beregningsmæssig kemi, fysisk motiveret modellering, samt algoritmiske metoder til søgning efter reaktionsveje i reaktionsnetværk. Resultatet er en udvidelse af et eksisterende software framework for graftransformation, som muliggør en mere realistisk udforskning af reaktionsnetværk. Søgningen efter reaktionsveje formuleres via mixed-integer linear programming (MILP), hvor kemiske potentialer og molekylkoncentrationer er tilføjet i optimeringsmodellen. Da molekylernes struktur er kendt, kan kemiske potentialer estimeres og integreres i optimeringsmodellen sammen med molekylkoncentrationer. Det gør det muligt at begrænse søgningen til reaktionsveje, der udelukkende består af termodynamisk favorable reaktioner. Når flere mulige reaktionsveje identificeres, kan de desuden rangordnes ved hjælp af en termodynamisk motiveret målfunktion.

Vi udvider dernæst metoden ovenfor ved at inddrage information om de enkelte reaktioners kinetik. De kinetiske parametre estimeres ud fra molekylernes struktur ved hjælp af en beregningsmæssig kemisk model. Da flere reaktionsveje ofte opfylder de angivne søgekriterier, introduceres en målfunktion baseret på fysiske principper om at maksimere den samlede sandsynlighed for en reaktionsvej. Det samlede resultat udgør en automatiseret analyse-pipeline, som tillader en fleksibel og kinetisk informeret analyse af reaktionsveje i store reaktionsnetværk, herunder reaktionsnetværk konstrueret via generative metoder hvor kinetiske data ikke medfølger.

Mange reaktionsnetværk er for store til, at de kan genereres og lagres eksplicit i deres helhed. For at håndtere denne udfordring introducerer denne afhandling et stokastisk beregningsmæssigt værktøj til udforskning af reaktionsnetværk, der er implicit defineret ved grafgrammatikker. Udforskningsstrategien ligner stokastisk kemisk kinetik ved at kombinere kollisionsbaseret udvælgelse af molekylepar med instantiering af reaktionsklasser. I hvert trin udvælger algoritmen først et molekylepar, der skal kollidere, derefter en reaktionsklasse og til sidst en konkret reaktionsinstans blandt de fundne muligheder for den valgte klasse. Denne kollisionsbaserede fremgangsmåde undgår at opliste alle mulige rektioner og gør en effektiv udforskning af store reaktionsnetværk under både åbne og lukkede systembetingelser mulig. Det udforskede delnetværk kan efterfølgende betragtes som et repræsentativt udsnit af det samlede molekylære netværk og kan analyseres ved hjælp af de metoder, der er udviklet i denne afhandling.

De foreslåede metoder anvendes derefter på en række kemiske eksempler, der belyser både styrker og begrænsninger for metoderne. Den termodynamisk informerede metode anvendes på et reaktionsnetværk, der modellerer hydrogencyanid-formamid kemi. Alternativer til de alment accepterede reaktionsveje identificeres og opregnes, herunder reaktionsveje som opnår en mere fordelagtig værdi af den termodynamiske målfunktion. Et reaktionsnetværk genereret ud fra glycolonitril, vand og ammoniak annoteres med beregnede aktiveringsenergier, hvorefter reaktionsveje til glycin og glykolsyre identificeres ved hjælp af den udviklede metode. Til sidst demonstreres det stokastiske værktøj til udforskning af store molekylære rum på  formose-kemi som eksempel. Dette studie undersøger både de kemiske egenskaber ved det udforskede delnetværk og de fordele, som kan opnås ved at anvende caching under udforskningen.

